\documentclass[20pt, a4paper]{article}
\usepackage[utf8]{inputenc}
\usepackage[T1]{fontenc}
\usepackage[italian]{babel}
\usepackage[table]{xcolor}
\usepackage{graphicx}
\usepackage{tikz}
\usetikzlibrary{
    shapes.geometric,
    arrows.meta,
    positioning,
    calc,
    patterns,
    decorations.pathmorphing
}
\usepackage{amsmath}
\usepackage{amssymb}
\usepackage{amsthm}
\usepackage{empheq}
\usepackage{siunitx}
\usepackage{textcomp}
\usepackage{array}
\usepackage{caption}
\usepackage{longtable}
\usepackage{multicol}
\usepackage{framed}
\usepackage{geometry}
\usepackage{enumitem}
\usepackage{tasks}
\usepackage{fancyhdr}
\geometry{
    a4paper,
    left=20mm,
    right=20mm,
    top=20mm,
    bottom=20mm,
}
\definecolor{headercolor}{RGB}{200,200,200}
\newcounter{quesito}
\setcounter{quesito}{0}
\newcommand{\nuovoquesito}{\stepcounter{quesito}\textbf{Domanda \arabic{quesito}.}}
\newcommand{\itemdomanda}[2]{%
    \item \begin{minipage}[t]{\linewidth}
        \nuovoquesito \\ #1
        \begin{enumerate}[label=(\Alph*), nosep, leftmargin=*]
            #2
        \end{enumerate}
    \end{minipage}
    \vspace{10pt}
}
\pagestyle{fancy}
\fancyhead[L]{Eserciziario}
\fancyhead[R]{\thepage}
\fancyfoot{}
\renewcommand{\headrulewidth}{0.4pt}

\begin{document}
\begin{titlepage}
    \centering
    \vspace*{2cm}
    {\Huge \textbf{Eserciziario TOLC-I} \par}
    \vspace{1cm}
    {\Large \textbf{Volume Matematica} \par}
    \vspace{4cm}
    {\Large A.A. 2025 -- 2026 \par}
    \vspace{4cm}
    \vfill
    {\large \copyright \ Scritto da Tutor \textbf{Matteo Zambito} \\
    per il progetto \textit{"Verso il Test di Ingegneria"} \par}
\end{titlepage}

\newpage
\tableofcontents
\newpage

\section{Matematica}
\begin{enumerate}[leftmargin=*]
\itemdomanda{$ (a-b)^3 $ è uguale a:
}{
\item $ a^3+3a^2b+3ab^2+b^3 $
\item $ a^3-b^3 $
\item $ a^3-3a^2b+3ab^2-b^3 $
\item $ a^3-3ab-b^3 $
\item $ a^3+3a^2b-3ab^2-b^3 $
}

\itemdomanda{La condizione di esistenza di $x^{-1} + y^{-1}$ è
}{
\item $ x \neq 1 \lor y \neq 1 $
\item $ x \neq -1 \lor y \neq -1 $
\item tutti i reali
\item $ x \neq y $
\item $ x \neq 0 \land y \neq 0 $
}

\itemdomanda{Qual è la condizione di esistenza dell'espressione $ \frac{2x}{x^2-4} $?
}{
\item $ x \neq 0 $
\item $ x \neq 2 $
\item $ x \neq 2 \land x \neq -2 $
\item $ x \neq 4 $
\item $ x \neq 0 \land x \neq 4 $
}

\itemdomanda{$(x-3)^{3} = $ ...
}{
\item $x^{3}\,-\,27$
\item $x^{3}\,-\,3x\,+\,27$
\item $x^{3}\,+\,3x\,+\,27$
\item $x^{3}\,-\,9x^{2}\,+\,27x\,-\,27$
\item $x^{3}\,-\,9x^{2}\,+\,27x\,+\,27$
}

\itemdomanda{Il risultato di $ (x+2)^3 $ è:
}{
\item $ x^3+8 $
\item $ x^3+6x^2+12x+6 $
\item $ x^3+6x^2+12x+8 $
\item $ x^3+4x^2+4x+8 $
\item $ x^3+6x^2+8 $
}

\itemdomanda{La frazione algebrica $ \frac{x^2 - 4x + 4}{x^2 - 4} $ è equivalente a:
}{
\item $ \frac{x+2}{x-2} $
\item $ \frac{x-2}{x+2} $
\item $ x+2 $
\item $ \frac{1}{x-2} $
\item $ 1 $
}

\itemdomanda{Se $ \frac{A}{B} = 0 $, allora:
}{
\item $ A=0 \land B=0 $
\item $ B=0 $
\item $ A=0 \land B \neq 0 $
\item $ A \neq 0 \land B=0 $
\item $ A \neq 0 $
}

\itemdomanda{L'espressione $ \frac{x^2-1}{x-1} $ per $ x \neq 1 $ è uguale a:
}{
\item $ x $
\item $ x+1 $
\item $ x-1 $
\item $ x^2 $
\item $ 1 $
}

\itemdomanda{Semplifica $ \frac{a^3b^2+a^2b^3}{a^2b^2} $ per $ a \neq 0, b \neq 0 $.
}{
\item $ ab $
\item $ a^2+b^2 $
\item $ a+b $
\item $ a^3+b^3 $
\item $ \frac{1}{ab} $
}

\itemdomanda{$\frac{2}{1-x} + \frac{2}{x-1}$ vale:
}{
\item $0$ per qualsiasi valore di $x$
\item $0$ solo se $x \neq 1$
\item $\frac{4}{1-x}$
\item $x \cdot \frac{4}{1-x}$
\item $\frac{4}{x-1}$
}

\itemdomanda{L'espressione $ \frac{1}{a} + \frac{1}{b} $ per $ a \neq 0, b \neq 0 $ è uguale a:
}{
\item $ \frac{2}{a+b} $
\item $ \frac{a+b}{ab} $
\item $ \frac{1}{ab} $
\item $ \frac{a+b}{a^2b^2} $
\item $ \frac{1}{a+b} $
}

\itemdomanda{$ \frac{x}{x-1} - 1 = $ ...
}{
\item $ \frac{x-1}{x-1} $
\item $ \frac{1}{x-1} $
\item $ \frac{x-2}{x-1} $
\item $ 0 $
\item $ -1 $
}

\itemdomanda{L'espressione $ \frac{x^2+x}{x^2-1} $ per $ x \neq 1, x \neq -1 $ è uguale a:
}{
\item $ \frac{x}{x-1} $
\item $ \frac{x^{2}}{x-1} $
\item $ \frac{x+1}{x-1} $
\item $ \frac{x}{x+1} $
\item $ 1 $
}

\itemdomanda{Qual è il grado complessivo del polinomio $P(x,y) = 5x^{4}y^{2} - 2x^{3}y^{5} + x^{6}$?}{
\item $8$
\item $6$
\item $15$
\item $4$
\item $7$
}

\itemdomanda{Il MCD tra $ 6a^2b $ e $ 9ab^3 $ è:
}{
\item $ 54a^3b^4 $
\item $ 18a^2b^3 $
\item $ 3ab $
\item $ 6ab $
\item $ 9a^2b^3 $
}

\itemdomanda{Il minimo comune multiplo tra $(a+b)$, $(a+2b)$, $(2a+b)$ e $(2a+2b)$ è
}{
\item $(2a+2b)$
\item $(2a+b)(a+2b)$
\item $(2a+2b)(a+b)$
\item $(a+b)(a+2b)(2a+b)$
\item $2(a+b)(a+2b)(2a+b)$
}

\itemdomanda{Il mcm tra $ 6a^2b $ e $ 9ab^3 $ è:
}{
\item $ 3ab $
\item $ 54a^3b^4 $
\item $ 18a^2b^3 $
\item $ 6a^2b^3 $
\item $ 9a^2b $
}

\itemdomanda{Il MCM tra $ x^2 $ e $ x+1 $ è:
}{
\item $ x $
\item $ 1 $
\item $ x^2(x+1) $
\item $ x^2+1 $
\item $ x(x+1) $
}

\itemdomanda{Il monomio $4abc^{3}$ è di:
}{
\item $1^{\circ}$ grado
\item $2^{\circ}$ grado
\item $3^{\circ}$ grado
\item $4^{\circ}$ grado
\item $5^{\circ}$ grado
}

\itemdomanda{Il grado del monomio $ -5x^2y^3z $ è:
}{
\item $ 2 $
\item $ 3 $
\item $ 6 $
\item $ 5 $
\item $ 1 $
}

\itemdomanda{Il coefficiente di $ x^2 $ nel polinomio $ (x^2+2x+1)(x-3) $ è:
}{
\item $ 2 $
\item $ 1 $
\item $ -1 $
\item $ -3 $
\item $ 0 $
}

\itemdomanda{Il polinomio $12x^{3}+2x^{2}+4x+14$ è divisibile per il binomio
}{
\item $x+1$
\item $x-1$
\item $x+2$
\item $x-2$
\item $x-3$
}

\itemdomanda{Quale dei seguenti binomi divide il polinomio $ P(x)=x^3-x^2-x+1 $?
}{
\item $ x-1 $
\item $ x+2 $
\item $ x+1 $
\item $ x-2 $
\item $ x $
}

\itemdomanda{Il polinomio $ P(x)=x^3-8 $ è divisibile per:
}{
\item $ x+2 $
\item $ x-2 $
\item $ x-8 $
\item $ x $
\item $ x+8 $
}

\itemdomanda{Il quoziente della divisione tra $ P(x)=x^3+x^2+x+1 $ e $ D(x)=x+1 $ è:
}{
\item $ x^2+1 $
\item $ x^2+2 $
\item $ x^2+x $
\item $ x^2+2x+1 $
\item $ x^2-1 $
}

\itemdomanda{Il grado del polinomio $ P(x,y)=x^3y^2+2x^4y-5xy^5 $ è:
}{
\item $ 4 $
\item $ 5 $
\item $ 6 $
\item $ 3 $
\item $ 2 $
}

\itemdomanda{Il resto della divisione di $ P(x)=x^2+x-2 $ per $ x+1 $ è:
}{
\item $ 0 $
\item $ 1 $
\item $ -2 $
\item $ 2 $
\item $ -1 $
}

\itemdomanda{Il monomio $(3x^{2}y^{3})^{2}$ è uguale a:}{
\item $9x^{4}y^{6}$
\item $6x^{4}y^{6}$
\item $9x^{4}y^{5}$
\item $3x^{4}y^{6}$
\item $9x^{2}y^{5}$
}

\itemdomanda{Il risultato di $ 9991 \cdot 10009 $ è:
}{
\item $99999119$
\item $99999999$
\item $99991119$
\item $99999919$
\item $91999919$
}

\itemdomanda{Il risultato di $ (1008)^{2} $ è
}{
\item $1016064$
\item $1016164$
\item $1016264$
\item $1016364$
\item $1016464$
}

\itemdomanda{$ (x-y)(x+y) = $ ...
}{
\item $ x^2+y^2 $
\item $ x^2-y^2 $
\item $ x^2-2xy+y^2 $
\item $ x^2+2xy+y^2 $
\item $ x^2-xy+y^2 $
}

\itemdomanda{$ (x+1)^2 - (x-1)^2 = $ ...
}{
\item $ 0 $
\item $ 2 $
\item $ 4x $
\item $ 2x^2+2 $
\item $ -4x $
}

\itemdomanda{Il risultato di $ (2x+1)^2 $ è:
}{
\item $ 4x^2+1 $
\item $ 4x^2+4x+1 $
\item $ 2x^2+2x+1 $
\item $ 4x^2+2x+1 $
\item $ 4x^2+4x $
}

\itemdomanda{Il quadrato di un binomio $ (A+B)^2 $ è:
}{
\item $ A^2+B^2 $
\item $ A^2+AB+B^2 $
\item $ A^2+2AB+B^2 $
\item $ A^2-2AB+B^2 $
\item $ (A+B)(A-B) $
}

\itemdomanda{Il risultato di $ (a+b+c)^2 $ è:
}{
\item $ a^2+b^2+c^2+2ab+2ac+2bc $
\item $ a^2+b^2+c^2 $
\item $ a^2+b^2+c^2+ab+ac+bc $
\item $ a^2+b^2+c^2+2abc $
\item $ a^2+b^2+c^2+4ab+4ac+4bc $
}

\itemdomanda{La scomposizione del polinomio $x^{2} - y^{2} + x - y$ è:}{
\item $(x-y)(x+y+1)$
\item $(x+y)(x-y+1)$
\item $(x-y)^{2}$
\item $(x-y)(x+y-1)$
\item $(x+y)(x+y+1)$
}

\itemdomanda{Scomponi $ 4x^2 - 9 $.
}{
\item $ (2x+3)^2 $
\item $ (2x-3)(2x+3) $
\item $ (4x-9)(4x+9) $
\item $ (2x-3)^2 $
\item $ (4x-3)(x+3) $
}

\itemdomanda{L'espressione $ x^4-16 $ si scompone in:
}{
\item $ (x^2-4)^2 $
\item $ (x^2+4)(x-2)(x+2) $
\item $ (x-2)^4 $
\item $ (x^2-4)(x^2+4) $
\item $ (x-4)(x+4) $
}

\itemdomanda{Scomponendo il polinomio $ x^2-5x+6 $, si ottiene:
}{
\item $ (x-6)(x+1) $
\item $ (x-2)(x-3) $
\item $ (x+2)(x+3) $
\item $ (x-1)(x-6) $
\item $ (x-3)(x+2) $
}

\itemdomanda{L'espressione $ \frac{x}{x-2} \cdot \frac{x^2-4}{x^2} $ per i valori ammissibili è:
}{
\item $ x+2 $
\item $ \frac{x+2}{x} $
\item $ \frac{1}{x} $
\item $ \frac{x-2}{x} $
\item $ x(x+2) $
}

\itemdomanda{Quale delle seguenti scomposizioni è corretta?
}{
\item $ a^3-b^3 = (a-b)(a^2-ab+b^2) $
\item $ a^3+b^3 = (a+b)(a^2-ab+b^2) $
\item $ a^3+b^3 = (a+b)^3 $
\item $ a^3-b^3 = (a-b)^3 $
\item $ a^3-b^3 = (a-b)(a^2+b^2) $
}

\itemdomanda{$\frac{4c^{3}-4a^{2}c}{c-a} = $ ...
}{
\item $4c(c-a)$
\item $4c^{2}(c+a)$
\item $4c(c+a^{2})$
\item $4c(c^{2}+a^{2})$
\item $4c(c+a)$
}

\itemdomanda{Qual è il resto della divisione del polinomio $P(x) = x^{3} - 4x + 1$ per $(x - 2)$?}{
\item $1$
\item $0$
\item $-1$
\item $17$
\item $3$
}

\itemdomanda{Quale delle seguenti disuguaglianze è falsa:
}{
\item $ 10^{100} < 10^{101} $
\item $ 100^{1000} < 1000^{100} $
\item $ 100^{100} < 101^{101}$
\item $ 100^{100} < 101^{100} $
\item $ 100^{100} < 100^{1000} $
}

\itemdomanda{Quale delle seguenti disuguaglianze è vera?
}{
\item $ \frac{3}{4} < \frac{2}{3} $
\item $ \frac{5}{6} > 1 $
\item $ \frac{1}{5} < \frac{1}{4} $
\item $ 0,1 > \frac{1}{9} $
\item $ \frac{1}{2} > 0,5 $
}

\itemdomanda{Determinare la terza parte delle somme di $ 5/6 $ e $ 9/8 $:
}{
\item $ \frac{47}{24} $
\item $ \frac{47}{72} $
\item $ {47} $
\item $ \frac{24}{47} $
\item $ {24} $
}

\itemdomanda{$ \frac{2}{5} $ del numero $ 50 $ è:
}{
\item $ 10 $
\item $ 5 $
\item $ 20 $
\item $ 25 $
\item $ 40 $
}

\itemdomanda{Il risultato dell'operazione $ \frac{1}{2} - \frac{1}{3} + \frac{1}{6} $ è:
}{
\item $ 0 $
\item $ 1 $
\item $ \frac{1}{3} $
\item $ \frac{1}{6} $
\item $ \frac{5}{6} $
}

\itemdomanda{La frazione generatrice del numero decimale periodico $ 0,\bar{3} $ è:
}{
\item $ \frac{3}{10} $
\item $ \frac{1}{30} $
\item $ \frac{1}{3} $
\item $ \frac{3}{100} $
\item $ \frac{3}{99} $
}

\itemdomanda{Quale tra i seguenti numeri è compreso tra $ 0,2 $ e $ 0,3 $?
}{
\item $ \frac{1}{2} $
\item $ \frac{3}{10} $
\item $ \frac{1}{4} $
\item $ \frac{1}{5} $
\item $ \frac{1}{3} $
}

\itemdomanda{Quale delle seguenti affermazioni è vera?
}{
\item $ \frac{1}{3} > 0,334 $
\item $ 0,66 < \frac{2}{3} $
\item $ \frac{1}{4} < 0,2 $
\item $ \frac{3}{2} < 1,5 $
\item $ 0,5 = \frac{1}{5} $
}

\itemdomanda{Il reciproco di $ 1,2 $ è:
}{
\item $ -1,2 $
\item $ \frac{12}{10} $
\item $ \frac{5}{6} $
\item $ \frac{6}{5} $
\item $ 0,8 $
}

\itemdomanda{Il MCD (Massimo Comune Divisore) tra $ 12, 18, 30 $ è:
}{
\item $ 60 $
\item $ 12 $
\item $ 6 $
\item $ 2 $
\item $ 3 $
}

\itemdomanda{Il mcm (minimo comune multiplo) tra $ 12, 18, 30 $ è:
}{
\item $ 60 $
\item $ 180 $
\item $ 240 $
\item $ 360 $
\item $ 90 $
}

\itemdomanda{Il numero $ 0,000000125 $ in notazione scientifica è:
}{
\item $ 12,5 \cdot 10^{-8} $
\item $ 1,25 \cdot 10^{-7} $
\item $ 1,25 \cdot 10^{7} $
\item $ 12,5 \cdot 10^{-9} $
\item $ 1,25 \cdot 10^{-8} $
}

\itemdomanda{La radice quadrata del quadrato della radice quadrata del quadrato di $-4$:
}{
\item non esiste
\item è $-4$
\item è $4$
\item è $2$
\item è $16$
}

\itemdomanda{Qual è la percentuale che corrisponde alla frazione $ \frac{3}{8} $?
}{
\item $ 30\% $
\item $ 37\% $
\item $ 37,5\% $
\item $ 38\% $
\item $ 80\% $
}

\itemdomanda{Determinare il $15\%$ dell' $80\%$ di $200$:}{
\item $24$
\item $30$
\item $16$
\item $48$
\item $120$
}

\itemdomanda{Qual è il risultato dell'operazione: $ \left( 3^{-2} \right) \cdot \left( \frac{2}{3} \right)^{2} : \left( \frac{2}{3} \right)^{-2} $
}{
\item $ {\frac{2^{4}}{3^{6}}} $
\item $ {\frac{2^{4}}{3^{4}}} $
\item $ {\frac{2^{2}}{3^{4}}} $
\item $ {\frac{2}{3^{4}}} $
\item $ {\frac{2^{3}}{3^{4}}} $
}

\itemdomanda{L'espressione $ \frac{1}{2} \cdot 4^{-1} $ è uguale a:
}{
\item $ 2 $
\item $ \frac{1}{8} $
\item $ 8 $
\item $ \frac{1}{2} $
\item $ \frac{1}{4} $
}

\itemdomanda{L'espressione $ \frac{3^5 \cdot 9^2}{3^3} $ è uguale a:
}{
\item $ 3^4 $
\item $ 3^{10} $
\item $ 3^6 $
\item $ 9^2 $
\item $ 3^5 $
}

\itemdomanda{$ 10^{-2} \cdot 10^3 \cdot 10^{-1} $ è uguale a:
}{
\item $ 10^{-6} $
\item $ 10 $
\item $ 1 $
\item $ 10^0 $
\item $ 10^2 $
}

\itemdomanda{Il valore di $ \left( \frac{1}{4} \right)^{-2} $ è:
}{
\item $ \frac{1}{16} $
\item $ -16 $
\item $ 16 $
\item $ \frac{1}{8} $
\item $ 8 $
}

\itemdomanda{Se $ x = 2^{-1} $, quanto vale $ x^2 $?
}{
\item $ -4 $
\item $ 4 $
\item $ \frac{1}{4} $
\item $ \frac{1}{2} $
\item $ -\frac{1}{4} $
}

\itemdomanda{L'espressione $ 2^3 + 2^4 $ è uguale a:
}{
\item $ 2^7 $
\item $ 24 $
\item $ 2^{12} $
\item $ 32 $
\item $ 16 $
}

\itemdomanda{Qual è il valore dell'espressione $25^{\frac{1}{2}} \cdot 8^{\frac{1}{3}}$?}{
\item $10$
\item $13$
\item $20$
\item $40$
\item $5$
}

\itemdomanda{$ [(\frac{1}{2} + \frac{1}{2})^{-1} \,-\,1 ]^{100} $ equivale a:
}{
\item $ 0 $
\item $ 1 $
\item $ -1 $
\item $ 100 $
\item $ -100 $
}

\itemdomanda{$ \sqrt{39} $ è compreso tra
}{
\item $ 5 $ e $ 6 $
\item $ 6 $ e $ 7 $
\item $ 14 $ e $ 15 $
\item $ 18 $ e $ 19 $
\item $ 19 $ e $ 20 $
}

\itemdomanda{$ \sqrt{200} + \sqrt{2} = $ ...
}{
\item $ \sqrt{400} $
\item $ 11 \sqrt{2} $
\item $ \sqrt{202} $
\item $ \sqrt{220} $
\item $ 2\sqrt{10}+\sqrt{2} $
}

\itemdomanda{$ \sqrt{\sqrt{3}} = $ ...
}{
\item $ \sqrt{3} $
\item $ \sqrt{6} $
\item $ \sqrt[4]{3} $
\item $ \sqrt[4]{6} $
\item $ \sqrt[3]{3} $
}

\itemdomanda{$ \sqrt{0,09} = $ ...
}{
\item $ 0,003 $
\item $ 0,03 $
\item $ 0,3 $
\item $ 3^{-1} $
\item $ 3^{-2} $
}

\itemdomanda{Il valore di $ \sqrt[3]{\frac{8}{27}} $ è:
}{
\item $ \frac{4}{9} $
\item $ \frac{2}{3} $
\item $ \frac{8}{3} $
\item $ \frac{2}{9} $
\item $ \frac{4}{3} $
}

\itemdomanda{Qual è il risultato di $ \sqrt{100 - 64} $?
}{
\item $ 6 $
\item $ 18 $
\item $ 64 $
\item $ 36 $
\item $ 8 $
}

\itemdomanda{$ \sqrt{72} $ può essere scritto come:
}{
\item $ 8\sqrt{9} $
\item $ 9\sqrt{8} $
\item $ 6\sqrt{2} $
\item $ 2\sqrt{6} $
\item $ 3\sqrt{8} $
}

\itemdomanda{$ \sqrt[4]{16} $ è uguale a:
}{
\item $ 4 $
\item $ 8 $
\item $ 2 $
\item $ 1 $
\item $ \frac{1}{4} $
}

\itemdomanda{La radice cubica di $ 0,008 $ è:
}{
\item $ 0,02 $
\item $ 0,2 $
\item $ 0,002 $
\item $ 0,000008 $
\item $ 0,8 $
}

\itemdomanda{$ \sqrt[3]{a \cdot b} $ esiste
}{
\item solo per $a\geq 0$ e $b \leq 0$
\item solo per $a \leq 0$ e $b \geq 0$
\item solo per $a \cdot b$ positivo
\item solo per $a \cdot b$ negativo
\item sempre
}

\itemdomanda{$ \sqrt[8]{a \cdot b} $ esiste
}{
\item solo per a positivo e b negativo
\item solo per $a < 0$ e $b > 0$
\item solo per a e b concordi o $a = 0$ e $b = 0$
\item solo per $a \cdot b \geq 0$
\item sempre
}

\itemdomanda{Dati i radicali $ \sqrt[4]{100} ; \sqrt[3]{1000} ; \sqrt{11} $ si ha
}{
\item $ \sqrt[4]{100} < \sqrt[3]{1000} < \sqrt{11} $
\item $ \sqrt[3]{1000} < \sqrt{11} < \sqrt[4]{100} $
\item $ \sqrt{11} < \sqrt[4]{100} < \sqrt[3]{1000} $
\item $ \sqrt[4]{100} < \sqrt[3]{1000} < \sqrt{11} < \sqrt[3]{1000} $
\item $ \sqrt[4]{100} < \sqrt{11} < \sqrt[3]{1000} $
}

\itemdomanda{Il valore di $ \left( \sqrt{5} - \sqrt{3} \right) \left( \sqrt{5} + \sqrt{3} \right) $ è:
}{
\item $ 8 $
\item $ 2 $
\item $ 5 - 3\sqrt{15} $
\item $ \sqrt{2} $
\item $ 8 - 2\sqrt{15} $
}

\itemdomanda{Semplificando $ \sqrt{\frac{1}{2}} $, si ottiene:
}{
\item $ \frac{1}{2} $
\item $ \frac{\sqrt{2}}{4} $
\item $ \frac{\sqrt{2}}{2} $
\item $ 2\sqrt{2} $
\item $ \sqrt{2} $
}

\itemdomanda{$ \frac{3}{\sqrt{3}} = $ ...
}{
\item $ 3 \cdot \frac{\sqrt{3}}{2} $
\item $ \sqrt{3}$
\item $ \frac{1}{\sqrt{3}}$
\item $ \frac{1}{\sqrt[3]{3}}$
\item $ \frac{1}{3}$
}

\itemdomanda{L'espressione $ \frac{1}{\sqrt{3} - 1} $ razionalizzata è:
}{
\item $ \frac{\sqrt{3}+1}{4} $
\item $ \frac{\sqrt{3}+1}{2} $
\item $ \frac{\sqrt{3}-1}{2} $
\item $ \sqrt{3}+1 $
\item $ \sqrt{3}-1 $
}

\itemdomanda{Qual è il risultato della somma $\frac{1}{4} + \frac{1}{5} + \frac{1}{20}$?}{
\item $\frac{1}{2}$
\item $\frac{3}{29}$
\item $\frac{1}{9}$
\item $1$
\item $\frac{7}{20}$
}

\itemdomanda{Nell'insieme dei numeri reali la disequazione $ \sqrt{\sqrt{|-x|}} < 1 $
}{
\item $ x < 1 $
\item $ x > -1 $
\item nessun valore di $ x \in \mathbb{R} $
\item $ x < - 1 \lor x > 1$
\item $ -1 < x < 1 $
}

\itemdomanda{La disequazione $ \sqrt{|x|} > 0 $ è verificata per
}{
\item $ x > 0 $
\item $ x > 1 $
\item $ 0 < x < 1 $
\item $ x < 0 \lor x > 1 $
\item $ x \neq 0 $
}

\itemdomanda{$ \sqrt{x} > \sqrt{y} $ è verificata
}{
\item per $ x > y $
\item per $ x > y $ solo se $ x > 1 $ e $ y \ge 1 $
\item per $ x > y $ solo se $ x > 0 $ e $ y \ge 0 $
\item mai
\item sempre
}

\itemdomanda{La disequazione $ \sqrt{x} \cdot \sqrt{-x} \cdot \sqrt{|x|} < 1 $ è soddisfatta per
}{
\item $ x < 1 $
\item $ 0 < x < 1 $
\item $ - 1 < x < 0 $
\item $ - 1 < x < 1 $
\item $ x = 0 $
}

\itemdomanda{Si può affermare che $ \sqrt{x} \cdot \sqrt{y} > 0 $
}{
\item per $ x > 0 \land y < 0 $
\item per $ x < 0 \land y < 0 $
\item per qualsiasi valore di $x$ e $y$
\item per $ x > 0 \land y > 0 $
\item non ha soluzioni reali
}

\itemdomanda{La soluzione della disequazione $ |x-3| < 1 $ è:
}{
\item $ x < 4 $
\item $ 2 < x < 4 $
\item $ 2 < x < 6 $
\item $ x > 2 $
\item $ x < 2 \lor x > 4 $
}

\itemdomanda{La disequazione $ (x^{2} + 1 ) ( x^{2} - 4 ) > 0 $
}{
\item $ x > 4 $
\item $ x > -1 $
\item $ 1 < x \lor x < -1$
\item $ -2 < x < 2 $
\item $ x > 2 \lor x < - 2 $
}

\itemdomanda{La disequazione $ x^2-x-2 > 0 $ è verificata per:
}{
\item $ -1 < x < 2 $
\item $ x > 2 $
\item $ x < -1 \lor x > 2 $
\item $ x < -2 \lor x > 1 $
\item $ x < 1 \lor x > 2 $
}

\itemdomanda{La disequazione $ x^2+1 < 0 $ ha soluzioni reali?
}{
\item Sì, $ x=\pm 1 $
\item Sì, $ x=0 $
\item No
\item Sì, $ x \in \mathbb{R} $
\item Sì, $ x < 0 $
}

\itemdomanda{La disequazione $ x^2+4x+4 \geq 0 $ è verificata per:
}{
\item $ x \geq -2 $
\item $ x \leq -2 $
\item $ \forall x \in \mathbb{R} $
\item $ x \neq -2 $
\item non ha soluzioni
}

\itemdomanda{La disequazione $ \frac{x-1}{x+2} \leq 0 $ è verificata per:
}{
\item $ x \leq 1 $
\item $ x > -2 $
\item $ -2 < x \leq 1 $
\item $ x < -2 \lor x \geq 1 $
\item $ -2 \leq x \leq 1 $
}

\itemdomanda{La disequazione $ \frac{1}{x} > 0 $ è verificata per:
}{
\item $ x \neq 0 $
\item $ x < 0 $
\item $ x > 0 $
\item $ \forall x \in \mathbb{R} $
\item $ x \geq 1 $
}

\itemdomanda{Per quali valori di $x$ è verificata la disequazione $\frac{1}{x} < 1$?}{
\item $x < 0 \lor x > 1$
\item $x > 1$
\item $0 < x < 1$
\item $x \neq 0$
\item $x < 1$
}

\itemdomanda{L'equazione $ x^4-5x^2+4=0 $ ha soluzioni reali:
}{
\item $ x=\pm 1 $
\item $ x=\pm 2 $
\item $ x=\pm 1, x=\pm 2 $
\item $ x=1, x=4 $
\item non ha soluzioni reali
}

\itemdomanda{Quali sono le soluzioni di $ x^{101} + 1 = 0$
}{
\item $\not\exists$ soluzioni reali
\item $1$
\item $-1$
\item $101$
\item $-101$
}

\itemdomanda{L'equazione $ x^3 = -27 $ ha soluzione:
}{
\item $ x=3 $
\item $ x=\pm 3 $
\item $ x=-3 $
\item $ x=9 $
\item non ha soluzioni reali
}

\itemdomanda{Quali sono le soluzioni dell'equazione $ x^{1010} + 1 = 0 $
}{
\item non esistono soluzioni reali
\item $1$
\item $-1$
\item $101$
\item $-101$
}

\itemdomanda{L'equazione $\sqrt[11]{-x^{11}} = 121$
}{
\item ha come soluzioni $+11$ e $-11$
\item ha come soluzione $+11$
\item ha come soluzione $-11$
\item ha come soluzione $+121$
\item ha come soluzione $-121$
}

\itemdomanda{Nell'insieme dei numeri reali l'equazione $ \sqrt{\sqrt{|-x|}} = 1 $ è verificata per
}{
\item $ x = 1 $
\item $ x = 1 \lor x = -1 $
\item $ x = 0 \lor x = -1 \lor x = 1 $
\item nessun valore reale di $x$
\item $ \not\exists \, x \in \mathbb{R} $
}

\itemdomanda{L'equazione $ |2x-1| = 3 $ ha soluzioni:
}{
\item $ x=2 $
\item $ x=-1 $
\item $ x=2 \lor x=-1 $
\item $ x=1 $
\item non ha soluzioni
}

\itemdomanda{L'equazione $|x + 5| = -2$ ha come soluzioni:}{
\item Nessuna soluzione reale
\item $x = -7$
\item $x = -3$
\item $x = -7 \lor x = -3$
\item $x = 0$
}

\itemdomanda{L'equazione $ x^2-5x+6=0 $ ha soluzioni:
}{
\item $ x=1, x=6 $
\item $ x=-2, x=3 $
\item $ x=2, x=3 $
\item $ x=2, x=-3 $
\item non ha soluzioni reali
}

\itemdomanda{L'equazione $ x^2+9=0 $ ha soluzioni reali?
}{
\item Sì, $ x=3 $ e $ x=-3 $
\item Sì, solo $ x=3 $
\item No
\item Sì, solo $ x=0 $
\item Sì, $ x=\pm \sqrt{3} $
}

\itemdomanda{L'equazione $ 3x^2-6x=0 $ ha soluzioni:
}{
\item $ x=0, x=-2 $
\item $ x=0, x=2 $
\item $ x=0, x=6 $
\item $ x=3, x=6 $
\item $ x=3, x=2 $
}

\itemdomanda{Quale equazione di secondo grado ha soluzioni $ x_1 = 1 $ e $ x_2 = -2 $?
}{
\item $ x^2+x-2=0 $
\item $ x^2+x-1=0 $
\item $ x^2-x-2=0 $
\item $ x^2-3x+2=0 $
\item $ x^2+3x+2=0 $
}

\itemdomanda{L'equazione $x^{2} - 7x + 10 = 0$ ha come soluzioni:}{
\item $x=2, x=5$
\item $x=-2, x=-5$
\item $x=1, x=10$
\item $x=2, x=-5$
\item $x=3, x=4$
}

\itemdomanda{Il discriminante $ \Delta $ dell'equazione $ 2x^2-x+3=0 $ è:
}{
\item $ 23 $
\item $ 25 $
\item $ -23 $
\item $ 1 $
\item $ 0 $
}

\itemdomanda{L'equazione $ \sqrt{x+1} = 2 $ ha soluzione:
}{
\item $ x=1 $
\item $ x=2 $
\item $ x=3 $
\item $ x=-1 $
\item $ x=0 $
}

\itemdomanda{L'equazione $ \frac{x-1}{x-1} = 1 $ è verificata per:
}{
\item $ \forall x \in \mathbb{R} $
\item $ x=1 $
\item $ x \neq 1 $
\item $ x=0 $
\item non ha soluzioni
}

\itemdomanda{Nel campo dei numeri reali, il numero degli zeri del polinomio $P(x) = x^{2} + a^{2}$ con $a \in \mathbb{R}$ è
}{
\item dipende da $a$ se $a \neq 0 $
\item pari a $2$ se $a = -1$
\item pari a $2$ se $a = 1$
\item pari a $2$ se $a$ è un quadrato perfetto
\item sempre pari a zero se $a \neq 0$
}

\itemdomanda{Il polinomio $ P(x)=x^2+4 $ ha zeri reali?
}{
\item Sì, $ x=2 $ e $ x=-2 $
\item Sì, solo $ x=2 $
\item No
\item Sì, solo $ x=-2 $
\item Sì, $ x=0 $
}

\itemdomanda{Il sistema di equazioni $ \begin{cases} x+y=5 \\ x-y=1 \end{cases} $ ha soluzione:
}{
\item $ x=2, y=3 $
\item $ x=3, y=2 $
\item $ x=4, y=1 $
\item $ x=1, y=4 $
\item non ha soluzioni
}

\itemdomanda{Il sistema $ \begin{cases} 2x-y=3 \\ 4x-2y=6 \end{cases} $ è:
}{
\item determinato
\item impossibile
\item indeterminato
\item omogeneo
\item di primo grado
}

\itemdomanda{Il sistema $ \begin{cases} x+2y=1 \\ x+2y=3 \end{cases} $ è:
}{
\item impossibile
\item determinato
\item indeterminato
\item omogeneo
\item di secondo grado
}

\itemdomanda{Il sistema di equazioni $\begin{cases} x + y = 10 \\ 2x - y = 2 \end{cases}$ ha come soluzione:}{
\item $x=4, y=6$
\item $x=6, y=4$
\item $x=2, y=8$
\item $x=5, y=5$
\item $x=3, y=7$
}

\itemdomanda{$y= \frac{x+1}{\sqrt{1 - x}}$ è :
}{
\item una funzione algebrica, razionale, fratta
\item una funzione algebrica, irrazionale
\item luna funzione algebrica, razionale, intera
\item una funzione trascendente
\item non è una funzione
}

\itemdomanda{$y= x + \ln(x)$ è :
}{
\item una funzione algebrica, razionale, fratta
\item una funzione algebrica, irrazionale
\item luna funzione algebrica, razionale, intera
\item una funzione trascendente
\item non è una funzione
}

\itemdomanda{$y= \frac{x ^{ 2}+1}{x}$ è :
}{
\item una funzione algebrica, razionale, fratta
\item una funzione algebrica, irrazionale
\item luna funzione algebrica, razionale, intera
\item una funzione trascendente
\item non è una funzione
}

\itemdomanda{Quale delle seguenti è una funzione trascendente?
}{
\item $ y = x^2+1 $
\item $ y = \sqrt{x} $
\item $ y = e^x $
\item $ y = \frac{x}{x-1} $
\item $ y = \frac{1}{x} $
}

\itemdomanda{Siano $ f(x) = x^{2} $ e $ g(x) = 2x $ due funzioni. La funzione composta $ f(g(x)) $ è:
}{
\item $ 2x^{2} $
\item $ 4x^{2} $
\item $ (4x)^{2} $
\item $ x^{2x} $
\item $ x^{2} + 2x $
}

\itemdomanda{Siano date le funzioni $ f(x)=x^{2} $ e $ g(x) = \sqrt{x} $. La funzione composta $ g(f(x)) $ è:
}{
\item $ x $
\item $ | x | $
\item $ x^{2} + \sqrt{x} $
\item $ x^{2} \cdot \sqrt{x} $
\item $ x^{2} - \sqrt{x} $
}

\itemdomanda{Se $ f(x) = x+3 $ e $ g(x) = 2x $, allora $ g(f(x)) $ è:
}{
\item $ 2x+3 $
\item $ 2x+6 $
\item $ x+6 $
\item $ 2x^2+6x $
\item $ x+5 $
}

\itemdomanda{Siano $f(x) = x + 1$ e $g(x) = x^{2}$. Quanto vale la funzione composta $g(f(2))$?}{
\item $9$
\item $5$
\item $6$
\item $4$
\item $3$
}

\itemdomanda{Il dominio della funzione $ f(x) = \frac{1}{x^2-1} $ è:
}{
\item $ x \neq 1 $
\item $ x \neq -1 $
\item $ x \neq 1 \land x \neq -1 $
\item $ x \in \mathbb{R} $
\item $ x > 1 $
}

\itemdomanda{Qual è il dominio della funzione $f(x) = \frac{1}{\sqrt{x}}$?}{
\item $x > 0$
\item $x \ge 0$
\item $x \neq 0$
\item Tutto $\mathbb{R}$
\item $x < 0$
}

\itemdomanda{Il dominio della funzione $ f(x) = \sqrt{\frac{x-3}{x}} $ è:
}{
\item $ x < 0 \lor x \geq 3 $
\item $ x \geq 3 $
\item $ x \leq 0 $
\item $ 0 < x \leq 3 $
\item $ x \in \mathbb{R} $
}

\itemdomanda{La funzione inversa di $ f(x) = 2x + 1 $ è:
}{
\item $ f^{-1}(x) = \frac{x + 1}{2} $
\item $ f^{-1}(x) = \frac{x - 1}{2} $
\item $ f^{-1}(x) = x + 2 $
\item $ f^{-1}(x) = x - 2 $
\item $ f^{-1}(x) = 2x - 1 $
}

\itemdomanda{La funzione inversa di $ f(x) = \frac{x}{2} $ è:
}{
\item $ f^{-1}(x) = \frac{2}{x} $
\item $ f^{-1}(x) = x+2 $
\item $ f^{-1}(x) = 2x $
\item $ f^{-1}(x) = x-2 $
\item $ f^{-1}(x) = x $
}

\itemdomanda{La funzione inversa di $y = e^{x}$ è:}{
\item $y = \ln(x)$
\item $y = e^{-x}$
\item $y = \frac{1}{e^{x}}$
\item $y = x^{e}$
\item $y = \log_{10}(x)$
}

\itemdomanda{La funzione $ f(x) = \frac{x}{x+1} $ è:
}{
\item pari
\item periodica
\item dispari
\item né pari né dispari
\item con dominio coincidente con tutto $\mathbb{R}$
}

\itemdomanda{La funzione $ f(x) = \frac{x}{\sin(x)} $ è:
}{
\item pari
\item periodica
\item dispari
\item né pari né dispari
\item con dominio coincidente con tutto $\mathbb{R}$
}

\itemdomanda{Quale funzione è dispari?
}{
\item $ f(x) = x^2 $
\item $ f(x) = |x| $
\item $ f(x) = x^3 $
\item $ f(x) = \cos(x) $
\item $ f(x) = x^2+1 $
}

\itemdomanda{La funzione $f(x) = x \cdot \sin(x)$ è:}{
\item Pari
\item Dispari
\item Né pari né dispari
\item Periodica di periodo $\pi$
\item Costante
}

\itemdomanda{La successione $ S_{n} = n^{2} + 1 , n \in \mathbb{N} $ è:
}{
\item crescente
\item decrescente
\item costante
\item né crescente né decrescente
\item decrescente in senso lato
}

\itemdomanda{Le variabili $x , y$ e $z$ sono variabili positive, legate l'una alle altre nel seguente modo:
	\begin{itemize}
	\item $x$ è direttamente proporzionale al cubo di $z$
	\item $y$ è direttamente proporzionale alla quarta potenza di $z$
	\end{itemize}
}{
\item Il cubo di $y$ è direttamente proporzionale alla quarta potenza di $x$
\item Il quadrato di $y$ è direttamente proporzionale alla quarta potenza di $x$
\item Il cubo di $y$ è direttamente proporzionale al cubo di $x$
\item Il cubo di $y$ è direttamente proporzionale al quadrato di $x$
\item Il quadrato di $y$ è direttamente proporzionale al quadrato di $x$
}

\itemdomanda{Se $ y $ è direttamente proporzionale a $ x $ e $ y=6 $ quando $ x=2 $, quale è la costante di proporzionalità?
}{
\item $ 2 $
\item $ 6 $
\item $ 3 $
\item $ 12 $
\item $ \frac{1}{3} $
}

\itemdomanda{Quale tra le seguenti funzioni gode della seguente proprietà: ``esistono due numeri $A$ e $B$ diversi e tali che $ f(A) = - f(B) $ '' ?
}{
\item $ f(x) = | x | $
\item $ f(x) = x^{2} + 1$
\item $ f(x) = \frac{1}{x^{4}} $
\item $ f(x) = x + 1$
\item $ f(x) = \frac{1}{x^{2}} $
}

\itemdomanda{Quale tra le seguenti funzioni gode della seguente proprietà: ``esistono due numeri $A$ e $B$ diversi e tali che $ f(A) = f(B) $ '' ?
}{
\item $ f(x) = x $
\item $ f(x) = x + 1$
\item $ f(x) = x^{3} $
\item $ f(x) = \frac{1}{x} $
\item $ f(x) = \frac{1}{x^{2}} $
}

\itemdomanda{$ y = | 1 - x | $ è :
}{
\item 
\[ y = \begin{cases}
    1 - x & \text{se } x \leq 1, \\
    x - 1 & \text{se } x > 1.
\end{cases} \]
\item 
\[ y = \begin{cases}
    1 - x & \text{se } x \geq 1, \\
    x - 1 & \text{se } x < 1.
\end{cases} \]
\item 
\[ y = \begin{cases}
    1 - x & \text{se } x \geq 0, \\
    x - 1 & \text{se } x < 0.
\end{cases} \]
\item 
\[ y = \begin{cases}
    x - 1 & \text{se } x \geq 0, \\
    1 - x & \text{se } x < 1.
\end{cases} \]
\item 
\[ y = \begin{cases}
    x & \text{se } x \geq 1, \\
    x - 1 & \text{se } x < 1.
\end{cases} \]
}

\itemdomanda{Nella funzione $y=\sin(x)$ la $y$ rappresenta:
}{
\item la funzione
\item la variabile indipendente
\item la variabile dipendente
\item il codominio
\item il dominio
}

\itemdomanda{Uno degli zeri della funzione $ f(x) = (x - k)(k^2+kx+x^7)$ , con $ k \in \mathbb{R} $ è sicuramente:
}{
\item $ k - 1 $
\item $ k + 1 $
\item $ k $
\item $ k^{7} $
\item $ k^{2} $
}

\itemdomanda{L'equazione della circonferenza con centro nell'origine e raggio $ r=4 $ è:
}{
\item $ x^2+y^2=2 $
\item $ x^2+y^2=4 $
\item $ x^2+y^2=16 $
\item $ (x-4)^2+y^2=16 $
\item $ x^2+y^2=8 $
}

\itemdomanda{La distanza tra i punti $ A(1, 0) $ e $ B(4, 4) $ è:
}{
\item $ 5 $
\item $ \sqrt{13} $
\item $ 3 $
\item $ 7 $
\item $ 25 $
}

\itemdomanda{Qual è l'equazione della retta passante per l'origine e perpendicolare alla retta $y = x$?}{
\item $y = -x$
\item $y = x$
\item $y = 0$
\item $x = 0$
\item $y = 2x$
}

\itemdomanda{Il punto medio del segmento di estremi $A(0, 0)$ e $B(4, 6)$ ha coordinate:}{
\item $(2, 3)$
\item $(4, 6)$
\item $(2, 2)$
\item $(1, 1.5)$
\item $(3, 2)$
}

\itemdomanda{L'equazione della retta passante per $ P(1, 2) $ e con pendenza $ m=3 $ è:
}{
\item $ y = 3x+2 $
\item $ y = 3x-1 $
\item $ y = 3x+5 $
\item $ y = 2x+3 $
\item $ y = 3x+1 $
}

\itemdomanda{La retta di equazione $ 2x-3y+6=0 $ ha intercetta $ y $ pari a:
}{
\item $ -2 $
\item $ 2 $
\item $ -3 $
\item $ 3 $
\item $ 0 $
}

\itemdomanda{La retta $ y=3x-5 $ interseca l'asse $ x $ nel punto di ascissa:
}{
\item $ x=5 $
\item $ x=-5 $
\item $ x=\frac{5}{3} $
\item $ x=3 $
\item $ x=-\frac{5}{3} $
}

\itemdomanda{La retta parallela all'asse $ x $ passante per $ P(1, 5) $ ha equazione:
}{
\item $ x = 1 $
\item $ y = 5 $
\item $ y = 1 $
\item $ x+y=6 $
\item $ y = x+4 $
}

\itemdomanda{La retta passante per $ A(1, 1) $ e $ B(3, 5) $ ha pendenza:
}{
\item $ 1 $
\item $ -2 $
\item $ 2 $
\item $ \frac{1}{2} $
\item $ 4 $
}

\itemdomanda{Il punto di intersezione delle rette $ y=x $ e $ y=2x-1 $ è:
}{
\item $ (0, 0) $
\item $ (1, 2) $
\item $ (1, 1) $
\item $ (2, 2) $
\item $ (2, 1) $
}

\itemdomanda{Una retta perpendicolare a $ y = \frac{3}{4}x+1 $ ha pendenza $ m $ uguale a:
}{
\item $ \frac{3}{4} $
\item $ -\frac{4}{3} $
\item $ \frac{4}{3} $
\item $ -\frac{3}{4} $
\item $ 1 $
}

\itemdomanda{Le rette di equazione $ y=2x+1 $ e $ y=-\frac{1}{2}x+3 $ sono:
}{
\item parallele
\item coincidenti
\item perpendicolari
\item incidenti (non perpendicolari)
\item nessuna delle precedenti
}

\itemdomanda{La funzione $ y = \arctan(x) $ ha codominio:
}{
\item $ \mathbb{R} $
\item $ [-\pi/2, \pi/2] $
\item $ (-\pi/2, \pi/2) $
\item $ [0, \pi] $
\item $ (-\infty, 0) $
}

\itemdomanda{Individuare l'unica affermazione vera:
}{
\item $ \cos(30^{\circ}) > \cos(60^{\circ}) $
\item $\sin(30^{\circ}) > \cos(30^{\circ}) $
\item $\sin(30^{\circ}) > \sin(60^{\circ}) $
\item $\cos(45^{\circ}) > \sin(45^{\circ}) $
\item $\cos(45^{\circ}) > \cos(30^{\circ}) $
}

\itemdomanda{Quanto vale $ \cot(1983\pi) $:
}{
\item $0$
\item $-1$
\item $1$
\item $3$
\item Non esiste
}

\itemdomanda{La disequazione $ \cos(x) > 1 $ è verificata per:
}{
\item $ x \in \mathbb{R} $
\item $ x > 0 $
\item Nessun $ x \in \mathbb{R} $
\item $ 0 < x < 2\pi $
\item $ x = 2k\pi, k \in \mathbb{Z} $
}

\itemdomanda{La funzione tangente di $x$ è definita per:
}{
\item $ x \neq \frac{\pi}{2} + k \pi , k \in \mathbb{Z} $
\item $ x \neq k \pi , k \in \mathbb{Z} $
\item $ x \neq \frac{\pi}{4} + k \pi , k \in \mathbb{Z} $
\item $ x \neq k \frac{\pi}{2} , k \in \mathbb{Z} $
\item $ x \neq \frac{\pi}{3} + k \pi , k \in \mathbb{Z} $
}

\itemdomanda{La funzione cotangente di $x$ è definita per:
}{
\item $ x \neq \frac{\pi}{2} + k \pi , k \in \mathbb{Z} $
\item $ x \neq k \pi , k \in \mathbb{Z} $
\item $ x \neq \frac{\pi}{4} + k \pi , k \in \mathbb{Z} $
\item $ x \neq k \frac{\pi}{2} , k \in \mathbb{Z} $
\item $ x \neq \frac{\pi}{3} + k \pi , k \in \mathbb{Z} $
}

\itemdomanda{La funzione cosecante di $x$ è definita per:
}{
\item $ x \neq \frac{\pi}{2} + k \pi , k \in \mathbb{Z} $
\item $ x \neq k \pi , k \in \mathbb{Z} $
\item $ x \neq \frac{\pi}{4} + k \pi , k \in \mathbb{Z} $
\item $ x \neq k \frac{\pi}{2} , k \in \mathbb{Z} $
\item $ x \neq \frac{\pi}{3} + k \pi , k \in \mathbb{Z} $
}

\itemdomanda{La funzione secante di $x$ è definita per:
}{
\item $ x \neq \frac{\pi}{2} + k \pi , k \in \mathbb{Z} $
\item $ x \neq k \pi , k \in \mathbb{Z} $
\item $ x \neq \frac{\pi}{4} + k \pi , k \in \mathbb{Z} $
\item $ x \neq k \frac{\pi}{2} , k \in \mathbb{Z} $
\item $ x \neq \frac{\pi}{3} + k \pi , k \in \mathbb{Z} $
}

\itemdomanda{L'espressione $\cos^{2}(x) - \sin^{2}(x)$ è equivalente a:}{
\item $\cos(2x)$
\item $1$
\item $\sin(2x)$
\item $0$
\item $2\cos(x)$
}

\itemdomanda{La funzione $ y = \cos(x) $ :
}{
\item ha un grafico che non interseca mai il grafico della funzione $ y = \frac{4}{5} $
\item è periodica di periodo $ \pi $
\item è periodica di periodo $ \frac{\pi}{2} $
\item è crescente per $ 0 < x < \frac{\pi}{2} $
\item è decrescente per $ 0 < x < \pi $
}

\itemdomanda{La funzione $ y = \cot(x) $ :
}{
\item ha un grafico che non interseca mai il grafico della funzione $ y = 1 $
\item è periodica di periodo $ \pi $
\item è periodica di periodo $ \frac{\pi}{2} $
\item è crescente per $ 0 < x < \frac{\pi}{2} $
\item passa per l'origine
}

\itemdomanda{La funzione $ y = \sin(x) $ è :
}{
\item illimitata superiormente
\item periodica di periodo $ \pi $
\item periodica di periodo $ \frac{\pi}{2} $
\item crescente per $ 0 < x < \frac{\pi}{2} $
\item decrescente per $ 0 < x < \frac{\pi}{2} $
}

\itemdomanda{La funzione $ y = \tan(x) $ :
}{
\item ha un grafico che non interseca mai il grafico della funzione $ y = 1 $
\item è periodica di periodo $ \pi $
\item è periodica di periodo $ \frac{\pi}{2} $
\item è crescente per $ 0 < x < \frac{\pi}{2} $
\item non passa per l'origine
}

\itemdomanda{Quanto vale $ \arccos(\tan(\frac{\pi}{4})) $:
}{
\item $0$
\item $ \frac{\pi}{3} $
\item $ \frac{\pi}{4} $
\item $ \pi $
\item $1$
}

\itemdomanda{Quanto vale $ \arcsin(\cos(\arcsin(\tan(\frac{\pi}{4})))) $:
}{
\item $0$
\item $ \frac{\pi}{3} $
\item $ \frac{\pi}{4} $
\item $ \pi $
\item $1$
}

\itemdomanda{Quanto vale $ \arcsin(\sin(\frac{\pi}{6})) $:
}{
\item Non esiste
\item $ \frac{\pi}{3} $
\item $ \frac{\pi}{6} $
\item $ \pi $
\item $1$
}

\itemdomanda{Quanto vale $ \arccos(\sin(\frac{\pi}{3})) $:
}{
\item Non esiste
\item $ \frac{\pi}{3} $
\item $ \frac{\pi}{6} $
\item $ \pi $
\item $1$
}

\itemdomanda{$ \sin^2(x) + \cos^2(x) = $ ...
}{
\item $ \sin(2x) $
\item $ 1 $
\item $ \cos(2x) $
\item $ 0 $
\item $ 2 $
}

\itemdomanda{Il valore massimo della funzione $ y = 3\sin(x) $ è:
}{
\item $ 1 $
\item $ 2 $
\item $ 3 $
\item $ \pi $
\item $ 3\pi $
}

\itemdomanda{La funzione $ y= \cos(x) $ è pari o dispari?
}{
\item Dispari
\item Pari
\item Nè pari nè dispari
\item Simmetrica rispetto all'origine
\item Dipende dall'intervallo
}

\itemdomanda{Il periodo della funzione $ y = \sin(2x) $ è:
}{
\item $ 2\pi $
\item $ \pi/2 $
\item $ \pi $
\item $ 4\pi $
\item $ 1 $
}

\itemdomanda{Quanto vale $ \cos(\pi) + \cos(2\pi) + \cos(3\pi) + \cos(4\pi) + \cos(5\pi) + \cos(6\pi) + \cos(7\pi) + \cos(8\pi) + \cos(9\pi) $:
}{
\item $-1$
\item $0$
\item $1$
\item $2$
\item $10$
}

\itemdomanda{Quanto vale $ \sin(\pi) + \sin(2\pi) + \sin(3\pi) + \sin(4\pi) + \sin(5\pi) + \sin(6\pi) + \sin(7\pi) + \sin(8\pi) + \sin(9\pi) $:
}{
\item $-1$
\item $0$
\item $1$
\item $2$
\item $10$
}

\itemdomanda{Quanto vale $ \tan(1973\pi) $:
}{
\item $0$
\item $-1$
\item $1$
\item $3$
\item Non esiste
}

\itemdomanda{Quanto vale $\tan(45^{\circ}) + \cot(45^{\circ})$?}{
\item $2$
\item $1$
\item $\sqrt{2}$
\item $0$
\item Non esiste
}

\itemdomanda{Quanto vale $ \cos ( 360^{\circ} )^{\sin( 30^{\circ})} $:
}{
\item $1$
\item $\sqrt{2}$
\item $\frac{1}{2}$
\item Non è un numero reale
\item $-1$
}

\itemdomanda{Se $ y = \sin(45^{\circ}) $ e $ x = \sin(-45^{\circ}) $, allora quanto vale $ y - x $ ?
}{
\item $\sqrt{2}$
\item $0$
\item $2$
\item $\frac{\sqrt{2}}{2}$
\item $-\frac{\sqrt{2}}{2}$
}

\itemdomanda{Se $ y = \cos(45^{\circ}) $ e $ x = \cos(-45^{\circ}) $, allora quanto vale $ y - x $ ?
}{
\item $0$
\item $\sqrt{2}$
\item $2$
\item $\frac{\sqrt{2}}{2}$
\item $-\frac{\sqrt{2}}{2}$
}

\itemdomanda{L'espressione $ \sin(\pi/4) $ vale:
}{
\item $ 1 $
\item $ \frac{\sqrt{2}}{2} $
\item $ \frac{1}{2} $
\item $ \frac{\sqrt{3}}{2} $
\item $ 0 $
}

\itemdomanda{L'espressione $ \tan(\pi/3) $ vale:
}{
\item $ 1 $
\item $ \frac{\sqrt{3}}{3} $
\item $ \sqrt{3} $
\item $ \frac{1}{2} $
\item $ 0 $
}

\itemdomanda{Quanto vale $\cot(120^\circ) $:
}{
\item $\sqrt{3}$
\item $ \frac{1}{\sqrt{3}} $
\item $ - \frac{\sqrt{3}}{3} $
\item $ -\sqrt{3} $
\item Non esiste
}

\itemdomanda{Nell'intervallo $[0, \pi/2]$, per quale valore di $x$ si ha $\sin(x) = \frac{1}{2}$?}{
\item $\frac{\pi}{6}$
\item $\frac{\pi}{3}$
\item $\frac{\pi}{4}$
\item $0$
\item $\frac{\pi}{2}$
}

\itemdomanda{Quanto vale $ ( 0.5 \cdot \cos(60^{\circ}))^{\sin(30^{\circ})} $:
}{
\item $1$
\item $\sqrt{2}$
\item $\frac{1}{2}$
\item Non è un numero reale
\item $-1$
}

\itemdomanda{Quanto vale $ \cos ( 180^{\circ} )^{\sin( 30^{\circ})} $:
}{
\item $1$
\item $\sqrt{2}$
\item $\frac{1}{2}$
\item Non è un numero reale
\item $-1$
}

\itemdomanda{Per quali valori di $ x \in \mathbb{R} $ è verificata la disequazione $ 2 ^{ 2x+3} > 0 $
}{
\item $ x > - \frac{3}{2} $
\item $ x < - \frac{3}{2} $
\item $ x > \frac{3}{2} $
\item $ x < \frac{3}{2} $
\item $ \forall x \in \mathbb{R} $
}

\itemdomanda{La disequazione $ e^x > 1 $ è verificata per:
}{
\item $ x \in \mathbb{R} $
\item $ x < 0 $
\item $ x > 0 $
\item $ x > e $
\item $ x < e $
}

\itemdomanda{La disequazione $ \log_2(x) > 3 $ è verificata per:
}{
\item $ x > 6 $
\item $ x > 3 $
\item $ x > 8 $
\item $ 0 < x < 8 $
\item $ x < 8 $
}

\itemdomanda{L'equazione $ 2^{-x} = - 4 $ ha soluzione
}{
\item $ x = 2 $
\item $ x = -2 $
\item $ x = \frac{1}{2} $
\item $ x = - \frac{1}{2} $
\item non ha soluzioni reali
}

\itemdomanda{La relazione $ 2^{x} = 3^{x} $ è vera per
}{
\item nessun $ x \in \mathbb{R} $
\item $ x = \frac{3}{2} $
\item $ x = \frac{2}{3} $
\item $ x = 0 $
\item $ x = 1 $
}

\itemdomanda{$ 2^{x+1} = 16 $ ha soluzione:
}{
\item $ x=3 $
\item $ x=8 $
\item $ x=4 $
\item $ x=15 $
\item $ x=-1 $
}

\itemdomanda{Se $ 4^{x-1} = 1 $, allora $ x $ vale:
}{
\item $ 4 $
\item $ 0 $
\item $ 1 $
\item $ 2 $
\item $ -1 $
}

\itemdomanda{L'equazione $9^{x} = 27$ ha soluzione:}{
\item $x = \frac{3}{2}$
\item $x = 3$
\item $x = 2$
\item $x = \frac{2}{3}$
\item $x = \frac{1}{3}$
}

\itemdomanda{L'equazione $ \log_{x}(1000) = 2 $ è verificata per
}{
\item $ x = 10 $
\item $ x = \sqrt{10} $
\item $ x = 10\sqrt{10} $
\item $ x = 100\sqrt{10} $
\item $ x = 1000\sqrt{10} $
}

\itemdomanda{L'equazione $ \log(x^2+1) = 0 $ ha soluzione:
}{
\item $ x=0 $
\item $ x=\pm 1 $
\item $ x=2 $
\item $ x=\pm 2 $
\item non ha soluzioni reali
}

\itemdomanda{L'equazione $ \log_3(x+1) = \log_3(x-1) + 1 $ ha soluzione:
}{
\item $ x = 2 $
\item $ x = \frac{1}{2} $
\item $ x = 1 $
\item $ x = 0 $
\item non ha soluzioni reali
}

\itemdomanda{Qual è il valore di $x$ nell'equazione $\log_{2}(x) + \log_{2}(x - 2) = 3$?}{
\item $4$
\item $2$
\item $-2$
\item $8$
\item $6$
}

\itemdomanda{$ \log_{2}(200) = $ vale
}{
\item $100$
\item $ 3 + \log_{2}(25) $
\item $ 4 + \log_{2}(25) $
\item $ 3 + \log_{3}(25) $
\item $ 50 $
}

\itemdomanda{Il logaritmo $ \log_{10}(1000) $ vale:
}{
\item $ 2 $
\item $ 3 $
\item $ 10 $
\item $ 100 $
\item $ 0 $
}

\itemdomanda{Se $ \log_3(x) = 2 $, allora $ x $ vale:
}{
\item $ 6 $
\item $ 2 $
\item $ 9 $
\item $ \frac{2}{3} $
\item $ \sqrt{3} $
}

\itemdomanda{$ \log_{10}(0.01) $ vale:
}{
\item $ 2 $
\item $ 1 $
\item $ -2 $
\item $ -1 $
\item $ 0.1 $
}

\itemdomanda{Quale di questi numeri: 1) $ \log_{2}(3) $ ; 2) $ - \log_{3}(3) $ e 3) $ \log_{3}(-3) $ è privo di significato:
}{
\item solo 1) e 2)
\item solo l'1)
\item solo il 2)
\item solo il 3)
\item tutti
}

\itemdomanda{Uno solo dei seguenti numeri è compreso tra $2$ e $3$
}{
\item $ \log_{2}(3) $
\item $ \log_{2}(5) $
\item $ \log_{2}(9) $
\item $ \log_{2}(10) $
\item $ \log_{2}(11) $
}

\itemdomanda{Il dominio della funzione $ f(x) = \sqrt{\log_2(x-1)} $ è:
}{
\item $ x > 1 $
\item $ x \geq 2 $
\item $ x > 2 $
\item $ x > 0 $
\item $ x \geq 1 $
}

\itemdomanda{$\log(x^{2}) $ è ...
}{
\item $ 2 \log(x) $
\item $ 2 \log(-x) $
\item $ 2 \log(|x|) $
\item $ \log(2x) $
\item $ (\log(|x|))^{2} $
}

\itemdomanda{La relazione $ \log_a(x \cdot y) = $ ... , per $ a>0, a \neq 1, x>0, y>0 $ è:
}{
\item $ \log_a(x) \cdot \log_a(y) $
\item $ \log_a(x) + \log_a(y) $
\item $ \log_a(x) - \log_a(y) $
\item $ \log_a(x+y) $
\item $ \log_{a^2}(xy) $
}

\itemdomanda{$ \log_{a}(\frac{x}{y}) = $ ... , per $ a>0, a \neq 1, x>0, y>0 $ è:
}{
\item $ \frac{\log_a(x)}{\log_a(y)} $
\item $ \log_a(x) \cdot \log_a(y) $
\item $ \log_a(x) - \log_a(y) $
\item $ \log_a(x-y) $
\item $ \log_{a^2}(x/y) $
}

\itemdomanda{L'espressione $\ln(e^{5} / e^{2})$ è uguale a:}{
\item $3$
\item $e^{3}$
\item $\frac{5}{2}$
\item $7$
\item $10$
}

\itemdomanda{Considerando la progressione $ 1, 5, 9, 13, 17, 21, \dots $, individuare l'unica affermazione CORRETTA.
}{
\item L'elemento mancante è $25$ e la progressione è geometrica di ragione $4$
\item L'elemento mancante è $25$ e la progressione è aritmetica di ragione $4$
\item L'elemento mancante è $27$ e la progressione è geometrica di ragione $4$
\item L'elemento mancante è $27$ e la progressione è aritmetica di ragione $4$
\item L'elemento mancante è $25$ e la progressione è decrescente
}

\itemdomanda{Considerando la progressione $ 2, 6, 18, 54, \dots $, individuare l'unica affermazione CORRETTA.
}{
\item L'elemento mancante è $152$ e la progressione è geometrica di ragione $3$
\item L'elemento mancante è $152$ e la progressione è aritmetica di ragione $3$
\item L'elemento mancante è $162$ e la progressione è geometrica di ragione $3$
\item L'elemento mancante è $162$ e la progressione è aritmetica di ragione $3$
\item L'elemento mancante è $162$ e la progressione cresce linearmente
}

\itemdomanda{Il quinto termine di una progressione geometrica che inizia con $ 3, 6, 12, \dots $ è:
}{
\item $ 30 $
\item $ 24 $
\item $ 48 $
\item $ 36 $
\item $ 96 $
}

\itemdomanda{In una progressione geometrica $2, -4, 8, -16, \dots$, la ragione $q$ è:}{
\item $-2$
\item $2$
\item $-6$
\item $4$
\item $1/2$
}

\itemdomanda{La somma di $ 4, 7, 10, 13, 16, 19, 22, 25, 28, 31 $ è:
}{
\item $350$
\item $300$
\item $175$
\item $150$
\item $125$
}

\itemdomanda{La somma di $ 2, 6, 18, 54, 162, 486, 1458, 4374, 13122, 39366 $ è:
}{
\item $59048$
\item $58048$
\item $59028$
\item $58038$
\item $56048$
}

\itemdomanda{Qual è la somma dei primi 100 numeri interi positivi ($1+2+3+...+100$)?}{
\item $5050$
\item $5000$
\item $5100$
\item $10100$
\item $4950$
}

\itemdomanda{Considerando i dati in tabella l'unica relazione possibile tra le grandezze $x$ e $y$ è di:
\renewcommand{\arraystretch}{1.5}
\begin{tabular}{|>{\columncolor{headercolor}}c|c|c|c|c|c|c|}
\hline
$x$ & 2 & 4 & 6 & 8 & 10 & 12 \\ \hline
$y$ & 1 & 4 & 7 & 10 & 13 & 16 \\ \hline
\end{tabular}

}{
\item proporzionalità quadratica
\item proporzionalità inversa
\item proporzionalità diretta
\item crescita lineare
\item nessuna delle altre risposte
}

\itemdomanda{Considerando i dati in tabella l'unica relazione possibile tra le grandezze $x$ e $y$ è di:
\renewcommand{\arraystretch}{1.5}
\begin{tabular}{|>{\columncolor{headercolor}}c|c|c|c|c|c|c|}
\hline
$x$ & 2 & 4 & 6 & 8 & 10 & 12 \\ \hline
$y$ & 101 & 202 & 303 & 404 & 505 & 606 \\ \hline
\end{tabular}

}{
\item proporzionalità quadratica
\item proporzionalità inversa
\item proporzionalità diretta
\item crescita lineare
\item proporzionalità cubica
}

\itemdomanda{Considerando i dati in tabella l'unica relazione possibile tra le grandezze $x$ e $y$ è di:
\renewcommand{\arraystretch}{1.5}
\begin{tabular}{|>{\columncolor{headercolor}}c|c|c|c|c|c|c|}
\hline
$x$ & 2 & 4 & 6 & 8 & 10 & 12 \\ \hline
$y$ & 50 & 25 & $\frac{50}{3}$ & $\frac{25}{2}$ & 10 & $\frac{25}{3}$ \\ \hline
\end{tabular}

}{
\item proporzionalità quadratica
\item proporzionalità inversa
\item proporzionalità diretta
\item crescita lineare
\item nessuna delle altre risposte
}

\itemdomanda{Considerando i dati in tabella l'unica relazione possibile tra le grandezze $x$ e $y$ è di:
\renewcommand{\arraystretch}{1.5}
\begin{tabular}{|>{\columncolor{headercolor}}c|c|c|c|c|c|c|}
\hline
$x$ & 2 & 4 & 6 & 8 & 10 & 12 \\ \hline
$y$ & 4 & 16 & 36 & 64 & 100 & 144 \\ \hline
\end{tabular}

}{
\item proporzionalità quadratica
\item proporzionalità inversa
\item proporzionalità diretta
\item crescita lineare
\item proporzionalità cubica
}
\end{enumerate}

\newpage
\section*{Soluzioni Matematica}
\begin{itemize}[leftmargin=*]
\item [Matematica] Domanda 1: 3
\item [Matematica] Domanda 2: 5
\item [Matematica] Domanda 3: 3
\item [Matematica] Domanda 4: 4
\item [Matematica] Domanda 5: 3
\item [Matematica] Domanda 6: 2
\item [Matematica] Domanda 7: 3
\item [Matematica] Domanda 8: 2
\item [Matematica] Domanda 9: 3
\item [Matematica] Domanda 10: 2
\item [Matematica] Domanda 11: 2
\item [Matematica] Domanda 12: 2
\item [Matematica] Domanda 13: 1
\item [Matematica] Domanda 14: 1
\item [Matematica] Domanda 15: 3
\item [Matematica] Domanda 16: 5
\item [Matematica] Domanda 17: 3
\item [Matematica] Domanda 18: 3
\item [Matematica] Domanda 19: 5
\item [Matematica] Domanda 20: 3
\item [Matematica] Domanda 21: 3
\item [Matematica] Domanda 22: 1
\item [Matematica] Domanda 23: 1
\item [Matematica] Domanda 24: 2
\item [Matematica] Domanda 25: 1
\item [Matematica] Domanda 26: 3
\item [Matematica] Domanda 27: 3
\item [Matematica] Domanda 28: 1
\item [Matematica] Domanda 29: 4
\item [Matematica] Domanda 30: 1
\item [Matematica] Domanda 31: 2
\item [Matematica] Domanda 32: 3
\item [Matematica] Domanda 33: 2
\item [Matematica] Domanda 34: 3
\item [Matematica] Domanda 35: 1
\item [Matematica] Domanda 36: 1
\item [Matematica] Domanda 37: 2
\item [Matematica] Domanda 38: 2
\item [Matematica] Domanda 39: 2
\item [Matematica] Domanda 40: 2
\item [Matematica] Domanda 41: 2
\item [Matematica] Domanda 42: 5
\item [Matematica] Domanda 43: 1
\item [Matematica] Domanda 44: 2
\item [Matematica] Domanda 45: 3
\item [Matematica] Domanda 46: 2
\item [Matematica] Domanda 47: 3
\item [Matematica] Domanda 48: 3
\item [Matematica] Domanda 49: 3
\item [Matematica] Domanda 50: 3
\item [Matematica] Domanda 51: 2
\item [Matematica] Domanda 52: 3
\item [Matematica] Domanda 53: 3
\item [Matematica] Domanda 54: 2
\item [Matematica] Domanda 55: 2
\item [Matematica] Domanda 56: 3
\item [Matematica] Domanda 57: 3
\item [Matematica] Domanda 58: 1
\item [Matematica] Domanda 59: 1
\item [Matematica] Domanda 60: 2
\item [Matematica] Domanda 61: 3
\item [Matematica] Domanda 62: 3
\item [Matematica] Domanda 63: 3
\item [Matematica] Domanda 64: 3
\item [Matematica] Domanda 65: 2
\item [Matematica] Domanda 66: 1
\item [Matematica] Domanda 67: 1
\item [Matematica] Domanda 68: 2
\item [Matematica] Domanda 69: 2
\item [Matematica] Domanda 70: 3
\item [Matematica] Domanda 71: 3
\item [Matematica] Domanda 72: 2
\item [Matematica] Domanda 73: 1
\item [Matematica] Domanda 74: 3
\item [Matematica] Domanda 75: 3
\item [Matematica] Domanda 76: 2
\item [Matematica] Domanda 77: 5
\item [Matematica] Domanda 78: 4
\item [Matematica] Domanda 79: 5
\item [Matematica] Domanda 80: 2
\item [Matematica] Domanda 81: 3
\item [Matematica] Domanda 82: 2
\item [Matematica] Domanda 83: 2
\item [Matematica] Domanda 84: 1
\item [Matematica] Domanda 85: 5
\item [Matematica] Domanda 86: 5
\item [Matematica] Domanda 87: 3
\item [Matematica] Domanda 88: 5
\item [Matematica] Domanda 89: 4
\item [Matematica] Domanda 90: 2
\item [Matematica] Domanda 91: 5
\item [Matematica] Domanda 92: 3
\item [Matematica] Domanda 93: 3
\item [Matematica] Domanda 94: 3
\item [Matematica] Domanda 95: 3
\item [Matematica] Domanda 96: 3
\item [Matematica] Domanda 97: 1
\item [Matematica] Domanda 98: 3
\item [Matematica] Domanda 99: 3
\item [Matematica] Domanda 100: 3
\item [Matematica] Domanda 101: 1
\item [Matematica] Domanda 102: 5
\item [Matematica] Domanda 103: 2
\item [Matematica] Domanda 104: 3
\item [Matematica] Domanda 105: 1
\item [Matematica] Domanda 106: 3
\item [Matematica] Domanda 107: 3
\item [Matematica] Domanda 108: 2
\item [Matematica] Domanda 109: 1
\item [Matematica] Domanda 110: 1
\item [Matematica] Domanda 111: 3
\item [Matematica] Domanda 112: 3
\item [Matematica] Domanda 113: 3
\item [Matematica] Domanda 114: 5
\item [Matematica] Domanda 115: 3
\item [Matematica] Domanda 116: 2
\item [Matematica] Domanda 117: 3
\item [Matematica] Domanda 118: 1
\item [Matematica] Domanda 119: 1
\item [Matematica] Domanda 120: 2
\item [Matematica] Domanda 121: 4
\item [Matematica] Domanda 122: 1
\item [Matematica] Domanda 123: 3
\item [Matematica] Domanda 124: 2
\item [Matematica] Domanda 125: 2
\item [Matematica] Domanda 126: 2
\item [Matematica] Domanda 127: 1
\item [Matematica] Domanda 128: 3
\item [Matematica] Domanda 129: 1
\item [Matematica] Domanda 130: 1
\item [Matematica] Domanda 131: 2
\item [Matematica] Domanda 132: 3
\item [Matematica] Domanda 133: 1
\item [Matematica] Domanda 134: 4
\item [Matematica] Domanda 135: 1
\item [Matematica] Domanda 136: 3
\item [Matematica] Domanda 137: 1
\item [Matematica] Domanda 138: 1
\item [Matematica] Domanda 139: 1
\item [Matematica] Domanda 140: 3
\item [Matematica] Domanda 141: 4
\item [Matematica] Domanda 142: 5
\item [Matematica] Domanda 143: 1
\item [Matematica] Domanda 144: 3
\item [Matematica] Domanda 145: 3
\item [Matematica] Domanda 146: 3
\item [Matematica] Domanda 147: 1
\item [Matematica] Domanda 148: 1
\item [Matematica] Domanda 149: 1
\item [Matematica] Domanda 150: 2
\item [Matematica] Domanda 151: 2
\item [Matematica] Domanda 152: 3
\item [Matematica] Domanda 153: 2
\item [Matematica] Domanda 154: 3
\item [Matematica] Domanda 155: 3
\item [Matematica] Domanda 156: 2
\item [Matematica] Domanda 157: 3
\item [Matematica] Domanda 158: 3
\item [Matematica] Domanda 159: 1
\item [Matematica] Domanda 160: 5
\item [Matematica] Domanda 161: 3
\item [Matematica] Domanda 162: 1
\item [Matematica] Domanda 163: 2
\item [Matematica] Domanda 164: 2
\item [Matematica] Domanda 165: 1
\item [Matematica] Domanda 166: 1
\item [Matematica] Domanda 167: 5
\item [Matematica] Domanda 168: 2
\item [Matematica] Domanda 169: 4
\item [Matematica] Domanda 170: 2
\item [Matematica] Domanda 171: 1
\item [Matematica] Domanda 172: 1
\item [Matematica] Domanda 173: 3
\item [Matematica] Domanda 174: 3
\item [Matematica] Domanda 175: 2
\item [Matematica] Domanda 176: 3
\item [Matematica] Domanda 177: 2
\item [Matematica] Domanda 178: 3
\item [Matematica] Domanda 179: 1
\item [Matematica] Domanda 180: 2
\item [Matematica] Domanda 181: 1
\item [Matematica] Domanda 182: 1
\item [Matematica] Domanda 183: 1
\item [Matematica] Domanda 184: 1
\item [Matematica] Domanda 185: 1
\item [Matematica] Domanda 186: 2
\item [Matematica] Domanda 187: 3
\item [Matematica] Domanda 188: 3
\item [Matematica] Domanda 189: 1
\item [Matematica] Domanda 190: 3
\item [Matematica] Domanda 191: 4
\item [Matematica] Domanda 192: 5
\item [Matematica] Domanda 193: 3
\item [Matematica] Domanda 194: 3
\item [Matematica] Domanda 195: 5
\item [Matematica] Domanda 196: 4
\item [Matematica] Domanda 197: 1
\item [Matematica] Domanda 198: 3
\item [Matematica] Domanda 199: 1
\item [Matematica] Domanda 200: 3
\item [Matematica] Domanda 201: 1
\item [Matematica] Domanda 202: 1
\item [Matematica] Domanda 203: 1
\item [Matematica] Domanda 204: 2
\item [Matematica] Domanda 205: 2
\item [Matematica] Domanda 206: 3
\item [Matematica] Domanda 207: 3
\item [Matematica] Domanda 208: 4
\item [Matematica] Domanda 209: 2
\item [Matematica] Domanda 210: 2
\item [Matematica] Domanda 211: 3
\item [Matematica] Domanda 212: 2
\item [Matematica] Domanda 213: 3
\item [Matematica] Domanda 214: 1
\item [Matematica] Domanda 215: 2
\item [Matematica] Domanda 216: 3
\item [Matematica] Domanda 217: 3
\item [Matematica] Domanda 218: 1
\item [Matematica] Domanda 219: 3
\item [Matematica] Domanda 220: 1
\item [Matematica] Domanda 221: 1
\item [Matematica] Domanda 222: 4
\item [Matematica] Domanda 223: 3
\item [Matematica] Domanda 224: 2
\item [Matematica] Domanda 225: 1
\end{itemize}

\end{document}
